\section*{अध्याय \ref{ch:b1:det} --- सारणिक और रैखिक निकाय}

\begin{solution}{exo:b1:det:1}
$3 \times 2 - 1 \times 5 = 1$.

दूसरा: $L_2 \leftarrow L_2 - L_1$, $L_3 \leftarrow L_3 - L_2$ (मूल पंक्तियों
पर) पंक्तियाँ $(1,2,3), (3,3,3), (3,3,3)$ देते हैं: दो बराबर पंक्तियाँ, अतः
सारणिक $0$। (सारुस पुष्टि करता है: $45 + 84 + 96 - 105 - 48 - 72 = 0$।)

तीसरा: यह $x = 1, 2, 3$ वाला वांडरमोंड है (\cref{ex:b1:det:vandermonde}):
$(2-1)(3-1)(3-2) = 2$।
\end{solution}

\begin{solution}{exo:b1:det:2}
सारणिक (सभी स्तंभ पहले में जोड़कर, फिर गुणनखंडन करके) $(\lambda + 2)$ गुणा
है
\[
\begin{vmatrix}
1 & 1 & \lambda\\ 1 & \lambda & 1\\ 1 & 1 & 1
\end{vmatrix}
= -(\lambda - 1)^2
\]
($L_1 \leftarrow L_1 - L_3$, $L_2 \leftarrow L_2 - L_3$ से साफ़ कीजिए और
प्रसार कीजिए), जिससे $\det = -(\lambda+2)(\lambda-1)^2$ मिलता है। \omterm{def:b1:vspaces:free}{आधार} $\iff
\det \neq 0 \iff \lambda \notin \{1, -2\}$।
\end{solution}

\begin{solution}{exo:b1:det:3}
पहला निकाय: $\det = -7$; $x = \frac{1}{-7}\begin{vmatrix} 5 & 1\\ 4 &
-2\end{vmatrix} = \frac{-14}{-7} = 2$, $y = \frac{1}{-7}\begin{vmatrix} 2 &
5\\ 3 & 4\end{vmatrix} = \frac{-7}{-7} = 1$। जाँच: $2(2) + 1 = 5$; $3(2) - 2
= 4$।

दूसरा निकाय: $L_2 - L_1$ और $L_3 - 2L_1$ के बाद पंक्तियाँ $(1,1,1)$,
$(0,-2,0)$, $(0,-1,-3)$ हो जाती हैं, अतः
\[
\det A = \begin{vmatrix} 1&1&1\\ 1&-1&1\\ 2&1&-1\end{vmatrix}
= 1 \times \begin{vmatrix} -2 & 0\\ -1 & -3\end{vmatrix} = 6 .
\]
क्रामर, स्तंभों के स्थान पर $(6,2,1)^{\mathsf T}$ रखकर:
\[
x = \frac{6}{6} = 1, \qquad
y = \frac{12}{6} = 2, \qquad
z = \frac{18}{6} = 3
\]
(अंश उसी तरह संगणित)। जाँच: $1 + 2 + 3 = 6$; $1 - 2 + 3 = 2$; $2 + 2 - 3 =
1$।
\end{solution}

\begin{solution}{exo:b1:det:4}
संवर्धित आव्यूह का लघुकरण: $L_2 \leftarrow L_2 - 2L_1$, $L_3 \leftarrow L_3
- L_1$:
\[
\begin{pmatrix}
1 & 2 & -1 & 1 & 1\\
0 & 0 & 3 & -3 & 3\\
0 & 0 & 3 & -3 & 3
\end{pmatrix}
\to
\begin{pmatrix}
1 & 2 & -1 & 1 & 1\\
0 & 0 & 1 & -1 & 1\\
0 & 0 & 0 & 0 & 0
\end{pmatrix}.
\]
धुरी-अज्ञात $x, z$; मुक्त अज्ञात $y, t$। पश्च-प्रतिस्थापन: $z = 1 + t$, $x =
1 - 2y + z - t = 2 - 2y$। \omterm{def:b1:logic:sets}{हल-समुच्चय}:
\[
\{(2 - 2y,\; y,\; 1 + t,\; t) : y, t \in \R\}
= (2, 0, 1, 0) + \operatorname{Vect}\bigl((-2,1,0,0),\,
(0,0,1,1)\bigr),
\]
अर्थात् $\R^4$ का एक एफ़ीन समतल (विमा $2 = 4 - \operatorname{rk} 2$)।
\end{solution}

\begin{solution}{exo:b1:det:5}
आगमन; $n = 1$ रिक्त गुणनफल $= 1$ है। चरण के लिए $k = n, n-1, \dots, 2$ हेतु
$C_k \leftarrow C_k - x_1 C_{k-1}$ कीजिए (इसी क्रम में, ताकि हर संक्रिया अब
तक अनछुए स्तंभ का प्रयोग करे)। पहली पंक्ति $(1, 0, \dots, 0)$ हो जाती है; और
पंक्ति $i \geq 2$ में $k$-वीं प्रविष्टि $x_i^{k-1} - x_1 x_i^{k-2} =
x_i^{k-2}(x_i - x_1)$ हो जाती है। पहली पंक्ति के अनुदिश प्रसार करके और हर
पंक्ति $i$ से $(x_i - x_1)$ का गुणनखंडन करके:
\[
V(x_1, \dots, x_n)
= \prod_{i=2}^{n} (x_i - x_1)\cdot V(x_2, \dots, x_n),
\]
और आगमन-परिकल्पना गुणनफल $\prod_{i<j}(x_j - x_i)$ पूरा कर देती है।
\end{solution}

\begin{solution}{exo:b1:det:6}
$D_n$ का पहली पंक्ति के अनुदिश प्रसार: $D_n = 2 D_{n-1} -
1\cdot\begin{vmatrix} 1 & \ast\\ 0 & D_{n-2}\text{-खंड} \end{vmatrix}$; और
दूसरा सारणिक, अपने पहले स्तंभ के अनुदिश प्रसारित होकर, $D_{n-2}$ है। अतः
$D_n = 2D_{n-1} - D_{n-2}$, अर्थात् $D_n - D_{n-1} = D_{n-1} - D_{n-2}$:
अंतर अचर हैं और $D_2 - D_1 = 1$ के बराबर। इस प्रकार $D_n = D_1 + (n - 1) = n
+ 1$। (जाँच: $D_2 = 3$, और $3\times3$ वाली स्थिति \cref{ex:b1:det:cofactor}
है: $D_3 = 4$।)
\end{solution}

\begin{solution}{exo:b1:det:7}
$C_1 \leftarrow C_1 + C_2 + C_3$ पहले स्तंभ को अचर $(m+2)$ बना देता है; उसका
गुणनखंडन कीजिए:
\[
\det = (m+2)\begin{vmatrix}
1 & 1 & m\\ 1 & m & 1\\ 1 & 1 & 1
\end{vmatrix}
\overset{L_1 - L_3,\ L_2 - L_3}{=}
(m+2)\begin{vmatrix}
0 & 0 & m-1\\ 0 & m-1 & 0\\ 1 & 1 & 1
\end{vmatrix}
= (m+2)\cdot\bigl(-(m-1)^2\bigr)
\]
(पहले स्तंभ के अनुदिश प्रसार कीजिए: अकेली प्रविष्टि $1$ का चिह्न $+$ है, और
शेष $2 \times 2$ सारणिक $0 \cdot 0 - (m-1)(m-1) = -(m-1)^2$ है)।

विवेचन। $m \notin \{1, -2\}$: क्रामर निकाय; समीकरणों की सममिति से $x = y =
z$, और हर समीकरण $(m + 2)x = 1$ देता है: अद्वितीय हल $\bigl(\frac{1}{m+2},
\frac{1}{m+2}, \frac{1}{m+2}\bigr)$। $m = 1$: तीनों समीकरण एक ही बात $x + y
+ z = 1$ कहते हैं: अतः हल एफ़ीन समतल $x + y + z = 1$ बनाते हैं। $m = -2$:
तीनों समीकरण जोड़ने पर $0 = 3$ मिलता है: अतः \omterm{def:b1:logic:sets}{हल-समुच्चय} रिक्त।
\end{solution}

\begin{solution}{exo:b1:det:8}
($\Rightarrow$) यदि $A^{-1}$ की प्रविष्टियाँ पूर्णांक हैं: तो $\det A \cdot
\det A^{-1} = 1$, जहाँ दोनों सारणिक पूर्णांक हैं (प्रविष्टियों के गुणनफलों
के योग): और जिन दो पूर्णांकों का गुणनफल $1$ हो, वे दोनों $\pm1$ हैं।

($\Leftarrow$) सहगुणनखंड-सूत्र $A^{-1} = \frac{1}{\det
A}\operatorname{Com}(A)^{\mathsf T}$ ($n \leq 3$ के लिए सीधे प्रसार से जाँचा
गया, व्यापक रूप से मान लिया गया) में $\operatorname{Com}(A)$ की प्रविष्टियाँ
पूर्णांक हैं (हर सहगुणनखंड एक पूर्णांक सारणिक है); और $\det A = \pm 1$ से
भाग देने पर भी वे पूर्णांक बने रहते हैं।
\end{solution}

\begin{solution}{exo:b1:det:9}
सभी स्तंभ पहले में जोड़िए: नए पहले स्तंभ की हर प्रविष्टि $a + nb$ है; उसका
गुणनखंडन कीजिए, ताकि पहला स्तंभ पूरा इकाइयों का हो जाए। फिर \omterm{met:b1:matrices:gauss}{पंक्ति
संक्रियाएँ} $L_i \leftarrow L_i - L_1$ ($i \geq 2$) ऊपर-बाएँ $1$ के नीचे की
हर प्रविष्टि साफ़ कर देती हैं और उन पंक्तियों में विकर्ण पर $a$ तथा अन्यत्र
$0$ छोड़ देती हैं: अतः आव्यूह ऊपरी त्रिभुजाकार है, जिसका विकर्ण $(1, a,
\dots, a)$ है। इस प्रकार
\[
\det(aI + bJ) = (a + nb)\, a^{\,n-1} ,
\]
जो अशून्य है तभी जब $a \neq 0$ और $a + nb \neq 0$: अर्थात्
\cref{exo:b1:matrices:9} वाली शर्त।
\end{solution}

\begin{solution}{exo:b1:det:10}
खंड संक्रियाएँ (हर एक तदनुरूप $n$ अदिश संक्रियाओं का संयोजन, जिसकी अनुमति
\cref{thm:b1:det:props}~(1) देता है):
\[
\begin{vmatrix} A & B\\ B & A\end{vmatrix}
\overset{C_1 \leftarrow C_1 + C_2}{=}
\begin{vmatrix} A + B & B\\ A + B & A\end{vmatrix}
\overset{L_2 \leftarrow L_2 - L_1}{=}
\begin{vmatrix} A + B & B\\ 0 & A - B\end{vmatrix}
= \det(A+B)\,\det(A-B),
\]
अंतिम चरण के लिए खंड-त्रिभुजाकार नियम का प्रयोग करते हुए।
\end{solution}

\begin{solution}{exo:b1:det:11}
$C_1 \leftarrow C_1 + C_2 + C_3$ पहले स्तंभ को अचर $(a + b + c)$ बना देता
है; उसका गुणनखंडन कीजिए, फिर $L_2 \leftarrow L_2 - L_1$, $L_3 \leftarrow L_3
- L_1$:
\[
\Delta = (a+b+c)\begin{vmatrix}
1 & b & c\\
0 & a - b & b - c\\
0 & c - b & a - c
\end{vmatrix}
= (a+b+c)\bigl[(a-b)(a-c) + (b-c)^2\bigr],
\]
और प्रसार करने पर $(a-b)(a-c) + (b-c)^2 = a^2 + b^2 + c^2 - ab - bc - ca$।
$\C$ पर, $j^3 = 1$ और $1 + j + j^2 = 0$ के साथ:
\begin{align*}
(a + jb + j^2c)(a + j^2b + jc)
&= a^2 + b^2 + c^2 + (j + j^2)(ab + bc + ca)\\
&= a^2 + b^2 + c^2 - ab - bc - ca ,
\end{align*}
जिससे पूरा गुणनखंडन मिल जाता है। (संरचनात्मक रूप से: स्तंभ $(1, j,
j^2)^{\mathsf T}$ $M\,(1, j, j^2)^{\mathsf T} = (a + jb + j^2c)(1, j,
j^2)^{\mathsf T}$ संतुष्ट करता है, और इसी प्रकार $j^2$ तथा $1$ भी: तीनों
गुणनखंड उस चक्रीय आव्यूह के तीन ``अभिलक्षणिक मान'' हैं --- यह कथा स्नातक
वर्ष~2 के खंड में व्यवस्थित की गई है।)
\end{solution}

\begin{solution}{exo:b1:det:12}
$r = \operatorname{rk} A$ लिखिए।

\emph{कोई व्युत्क्रमणीय $r \times r$ उपआव्यूह है।} $A$ के $r$ \omterm{def:b1:vspaces:free}{स्वतंत्र} स्तंभ
चुनिए और मान लीजिए $B \in \mathcal{M}_{n,r}$ उनसे बना आव्यूह है:
$\operatorname{rk} B = r$। चूँकि पंक्ति-कोटि स्तंभ-कोटि के बराबर है
(\cref{thm:b1:matrices:rank}), अतः $B$ की $r$ पंक्तियाँ \omterm{def:b1:vspaces:free}{स्वतंत्र} हैं; उन्हीं
पंक्तियों को रखने पर $A$ का $r \times r$ उपआव्यूह मिलता है जिसकी कोटि $r$
है, अर्थात् वह व्युत्क्रमणीय है।

\emph{इससे बड़ा कोई नहीं।} मान लीजिए $S$ कोई व्युत्क्रमणीय $s \times s$
उपआव्यूह है, जो $A$ के स्तंभों $j_1, \dots, j_s$ और पंक्तियों $i_1, \dots,
i_s$ से लिया गया है। यदि तदनुरूप \emph{पूरे} स्तंभों का कोई संयोजन $\sum_k
\lambda_k C_{j_k} = 0$ शून्य हो जाए, तो केवल पंक्तियाँ $i_1, \dots, i_s$
पढ़ने पर $S$ के स्तंभों पर $\sum_k \lambda_k S_k = 0$ मिलता है, अतः सभी
$\lambda_k = 0$ ($S$ व्युत्क्रमणीय है): इस प्रकार $A$ के स्तंभ $C_{j_1},
\dots, C_{j_s}$ \omterm{def:b1:vspaces:free}{स्वतंत्र} हैं, और $s \leq \operatorname{rk} A = r$।

अतः $\operatorname{rk} A$ ठीक किसी व्युत्क्रमणीय उपआव्यूह का अधिकतम आकार है।
\end{solution}

\begin{solution}{pb:b1:det:1}
\textbf{1.} $V(1,2,3,4) = (2-1)(3-1)(4-1)(3-2)(4-2)(4-3) = 1 \cdot 2\cdot
3\cdot 1\cdot 2\cdot 1 = 12$। भिन्न गाँठों पर अंतर्वेशन $P$ के उन गुणांकों
की माँग करता है जो $W = (x_i^{\,j-1})$ के साथ $W c = y$ हल करें, और $\det W
= V \neq 0$: अर्थात् एक क्रामर निकाय।

\textbf{2.} स्तंभों पर बाएँ से दाएँ काम कीजिए। $C_1$ अचर स्तंभ $P_0(x_i) =
1$ है ($P_0$ $0$ घात का \omterm{def:b1:poly:def}{इकाई-अग्र} है)। मान लीजिए स्तंभ $1, \dots, j-1$ पहले
ही शुद्ध घातों $1, x_i, \dots, x_i^{\,j-2}$ तक घटाए जा चुके हैं। चूँकि
$P_{j-1} = X^{j-1} + \sum_{k < j-1}\alpha_k X^k$, अतः $C_j$ में से संयोजन
$\sum_k \alpha_k\,(\text{स्तंभ } x_i^k)$ घटाने पर --- यह संक्रिया सारणिक
नहीं बदलती --- शुद्ध घात-स्तंभ $x_i^{\,j-1}$ बच जाता है। अंतिम स्तंभ के बाद
आव्यूह वांडरमोंड आव्यूह है: $\det = V(x_1, \dots, x_n)$।

\textbf{3.} \omterm{def:b1:poly:def}{बहुपद} $(j-1)!\,B_{j-1}$ $j - 1$ घात के \omterm{def:b1:poly:def}{इकाई-अग्र} हैं, अतः प्रश्न
2 देता है
\[
\det\bigl(B_{j-1}(m_i)\bigr)
= \frac{V(m_1, \dots, m_n)}{0!\,1!\cdots(n-1)!} .
\]
बायाँ पक्ष \emph{पूर्णांक} प्रविष्टियों वाले किसी आव्यूह का सारणिक है ($B_k$
$\Z$ पर पूर्णांक-मान है: सप्ताहांत समस्या \cref{pb:b1:vspaces:1} के प्रश्न
16--17), अतः वह एक पूर्णांक है; और वह धनात्मक है, क्योंकि $m_1 < \dots <
m_n$ के लिए $V(m_1, \dots, m_n) > 0$। इस प्रकार अधिक्रमगुणित सभी युग्मशः
अंतरों के गुणनफल को \omterm{def:b1:arith:divides}{विभाजित करता है}।

\textbf{4.} हर पंक्ति $i$ से $x_i$ का गुणनखंडन कीजिए: $\det(x_i^{\,j})_{j =
1..n} = x_1\cdots x_n\, \det(x_i^{\,j-1}) = x_1\cdots x_n\,V$।

\textbf{5.} $(W^{\mathsf T}W)_{ij} = \sum_k x_k^{\,i-1} x_k^{\,j-1} =
p_{i+j-2}$: $S = W^{\mathsf T}W$। अतः $\det S = \det(W^{\mathsf T})\det W =
V^2$ (\cref{thm:b1:det:props}~(2),(4))। वास्तविक $x_i$ के लिए: $\det S = V^2
\geq 0$, और $S$ व्युत्क्रमणीय है तभी जब $V \neq 0$, अर्थात् तभी जब $x_i$
युग्मशः भिन्न हों --- यह एक चिह्न-निश्चित परीक्षा है जो अकेले घात-योगों से
संगणित हो जाती है।

\textbf{6.} अंतर्वेशन-शर्तें $\sum_{k} c_k\,x_i^{\,k} = y_i$ निकाय $Wc = y$
बनाती हैं; और $\det W = V \neq 0$ अस्तित्व तथा अद्वितीयता एक ही साथ दे देता
है। यह पुस्तक की तीसरी उपपत्ति है: \cref{thm:b1:poly:lagrange} में स्पष्ट
सूत्र, \cref{ex:b1:linmaps:interpolation} में अष्टि-तर्क, और यहाँ क्रामर।

\textbf{7.} क्रामर: $c_{n-1} = \det W'/\det W$, जहाँ $W'$ वही $W$ है जिसके
अंतिम स्तंभ के स्थान पर $y$ रख दिया गया है। उस स्तंभ के अनुदिश $\det W'$ का
प्रसार करने पर:
\[
\det W' = \sum_{i=1}^n (-1)^{i+n} y_i\,V(x_1, \dots, \widehat{x_i},
\dots, x_n) .
\]
अब $V = V(\setminus i)\cdot\prod_{j<i}(x_i - x_j)\prod_{j>i} (x_j - x_i)$,
और दूसरे गुणनफल को बदलने की क़ीमत $(-1)^{n-i}$ है:
\[
(-1)^{i+n}\,\frac{V(\setminus i)}{V}
= \frac{(-1)^{i+n}(-1)^{n-i}}{\prod_{j\neq i}(x_i - x_j)}
= \frac{1}{\prod_{j\neq i}(x_i - x_j)} ,
\]
जिससे $c_{n-1} = \sum_i y_i/\prod_{j \neq i}(x_i - x_j)$ --- अर्थात् फिर से
वही विभाजित-अंतर वाला सूत्र।

\textbf{8.} $L_3 \leftarrow L_3 - L_1$ पंक्तियाँ $(1, x_1, x_1^2)$, $(0, 1,
2x_1)$, $(0,\ x_2 - x_1,\ (x_2-x_1)(x_2+x_1))$ देता है; पहले स्तंभ के अनुदिश
प्रसार करके और $(x_2 - x_1)$ का गुणनखंडन करके:
\[
(x_2 - x_1)\begin{vmatrix} 1 & 2x_1\\ 1 & x_2 + x_1
\end{vmatrix} = (x_2 - x_1)(x_2 - x_1) = (x_2 - x_1)^2 .
\]
$x_1 \neq x_2$ के लिए यह अशून्य है: $P \in \R_2[X]$ के गुणांकों पर $P(x_1) =
u$, $P'(x_1) = v$, $P(x_2) = w$ को व्यक्त करने वाला \omterm{def:b1:det:system}{रैखिक निकाय} क्रामर है
--- अर्थात् दुहरी गाँठ वाला एरमीट अंतर्वेशन सुप्रस्तुत है।

\textbf{9.} $a = P(0) = 1$, $b = P'(0) = 0$, $a + b + c = P(1) = 2$ के साथ
$P = a + bX + cX^2$: $c = 1$, अतः $P = 1 + X^2$, और वह अद्वितीय है। संगति:
यहाँ $x_1 = 0$, $x_2 = 1$, और प्रश्न 8 का सारणिक $(1 - 0)^2 = 1 \neq 0$ है।

\textbf{10.} सीधी संगणना:
\[
C_2 = \frac{1}{(a_1+b_1)(a_2+b_2)} - \frac{1}{(a_1+b_2)(a_2+b_1)}
= \frac{(a_1+b_2)(a_2+b_1) - (a_1+b_1)(a_2+b_2)}
{\prod_{i,j}(a_i+b_j)} ,
\]
और अंश प्रसारित होकर $a_1b_1 + a_2b_2 - a_1b_2 - a_2b_1 = (a_2 - a_1)(b_2 -
b_1)$ हो जाता है: अर्थात् अंतर बटा योग।

\textbf{11.} $i < n$ के लिए पंक्ति $i$ की नई प्रविष्टि है
\[
\frac{1}{a_i + b_j} - \frac{1}{a_n + b_j}
= \frac{a_n - a_i}{(a_i + b_j)(a_n + b_j)} .
\]
हर पंक्ति $i < n$ से $(a_n - a_i)$ और फिर हर स्तंभ $j$ से $\frac1{a_n +
b_j}$ का गुणनखंडन कीजिए: जो बचता है उसकी पंक्तियों $i < n$ में प्रविष्टियाँ
$\frac1{a_i + b_j}$ हैं और पंक्ति $n$ में अचर $1$ --- अर्थात् आव्यूह $M$, और
साथ में वही घोषित पूर्व-गुणक।

\textbf{12.} $M$ पर, $j < n$ के लिए संक्रिया $C_j \leftarrow C_j - C_n$
पंक्ति $n$ को $(0, \dots, 0, 1)$ में बदल देती है, और पंक्ति $i < n$ में,
\[
\frac{1}{a_i + b_j} - \frac{1}{a_i + b_n}
= \frac{b_n - b_j}{(a_i + b_j)(a_i + b_n)} .
\]
हर स्तंभ $j < n$ से $(b_n - b_j)$ और हर पंक्ति $i < n$ से $\frac1{a_i +
b_n}$ का गुणनखंडन कीजिए, फिर अंतिम पंक्ति के अनुदिश प्रसार कीजिए (चिह्न
$(-1)^{n+n} = +1$): शेष सारणिक $C_{n-1}$ है। प्रश्न 11--12 के गुणनखंड इकट्ठे
करने पर:
\[
C_n = \frac{\prod_{i<n}(a_n - a_i)\,\prod_{j<n}(b_n - b_j)}
{\prod_{j}(a_n + b_j)\,\prod_{i<n}(a_i + b_n)}\;C_{n-1},
\]
और आगमन (\omterm{def:b1:vspaces:free}{आधार} $C_1 = \frac1{a_1+b_1}$) ठीक कोशी का दुहरा एकांतरक जोड़ देता
है: सभी जोड़ियों के लिए गुणनखंड $(a_j - a_i)(b_j - b_i)$, बटा सभी योग $(a_i
+ b_j)$।

\textbf{13.} सूत्र शून्य होता है तभी जब कोई $a_j = a_i$ अथवा $b_j = b_i$:
अर्थात् कोशी आव्यूह व्युत्क्रमणीय है तभी जब दोनों कुल युग्मशः भिन्न हों।
हिल्बर्ट: $a_i = i$, $b_j = j - 1$। $n = 2$ के लिए: अंश $(2-1)(1-0) = 1$, हर
$1\cdot2\cdot2\cdot3 = 12$: $\det H_2 = \frac1{12}$। $n = 3$ के लिए: अंश
$\bigl[(1)(2)(1)\bigr]^2 = 4$, हर $(1\cdot2\cdot3)
(2\cdot3\cdot4)(3\cdot4\cdot5) = 6\cdot24\cdot60 = 8640$: $\det H_3 =
\frac{4}{8640} = \frac1{2160}$। $n = 2$ के लिए प्रतिलोम:
\[
H_2^{-1} = 12\begin{pmatrix} \frac13 & -\frac12\\[2pt]
-\frac12 & 1\end{pmatrix}
= \begin{pmatrix} 4 & -6\\ -6 & 12 \end{pmatrix},
\]
सभी पूर्णांक (यह परिघटना हर $n$ के लिए सत्य है)।

\textbf{14.} निकाय का आव्यूह कोशी आव्यूह है, जो प्रश्न 13 से तब
व्युत्क्रमणीय है जब $b_j$ (और $a_i$) युग्मशः भिन्न हों: अतः अद्वितीय हल।
व्याख्या: सरल ध्रुवों वाला कोई परिमेय फलन $R = \sum_j \frac{c_j}{X + b_j}$
अपने $n$ मानों $R(a_1), \dots, R(a_n)$ से तय हो जाता है, और विलोमतः ऐसा हर
आँकड़ा-पत्रक ठीक एक बार साकार होता है --- यह आंशिक भिन्नों के अस्तित्व और
अद्वितीयता प्रमेय (\cref{thm:b1:fractions:complex}) का प्रतिचयन वाला समकक्ष
है।

\textbf{15.} $X^3 + pX + q$ के लिए वियेत: $\lambda_1 + \lambda_2 + \lambda_3
= 0$, $\sum_{i<j}\lambda_i\lambda_j = p$, अतः $p_1 = 0$ और $p_2 = p_1^2 - 2p
= -2p$। हर मूल $\lambda^3 = -p\lambda - q$ संतुष्ट करता है; योग करने पर:
$p_3 = -p\,p_1 - 3q = -3q$। $\lambda$ से गुणा करके योग करने पर: $p_4 =
-p\,p_2 - q\,p_1 = 2p^2$।

\textbf{16.} प्रश्न 5 से (सर्वसमिका $S = W^{\mathsf T}W$ और $\det S = V^2$
$\C$ पर मान्य हैं),
\[
V^2 = \begin{vmatrix}
3 & 0 & -2p\\
0 & -2p & -3q\\
-2p & -3q & 2p^2
\end{vmatrix}
= 3\bigl(-4p^3 - 9q^2\bigr) + (-2p)\bigl(0 - 4p^2\bigr)
= -4p^3 - 27q^2 ,
\]
पहली पंक्ति के अनुदिश प्रसार करते हुए।

\textbf{17.} दोहराए गए मूल का अर्थ है दो बराबर $\lambda_i$, अर्थात् $V = 0$,
यानी $\operatorname{disc} = -4p^3 - 27q^2 = 0$। $X^3 - 3X + 2$ के लिए:
$4(-3)^3 + 27\cdot4 = -108 + 108 = 0$, जो $(X-1)^2(X+2)$ के दुहरे मूल $1$ से
मेल खाता है।

\textbf{18.} किसी वास्तविक त्रिघातीय के अवास्तविक मूल संयुग्मी जोड़ियों में
आते हैं, अतः $\operatorname{disc} \neq 0$ होने पर ठीक दो स्थितियाँ बनती हैं।
तीन भिन्न वास्तविक मूल: $V$ वास्तविक और अशून्य है, अतः $\operatorname{disc}
= V^2 > 0$। एक वास्तविक मूल $\lambda_1$ तथा $\lambda_3 = \conj{\lambda_2}
\notin \R$: तब
\[
(\lambda_2 - \lambda_1)(\lambda_3 - \lambda_1) =
\abs{\lambda_2 - \lambda_1}^2 > 0,
\qquad
\lambda_3 - \lambda_2 = -2\iu\,\operatorname{Im}\lambda_2 \neq 0,
\]
अतः $V$ एक अशून्य विशुद्ध काल्पनिक संख्या है और $\operatorname{disc} = V^2 <
0$। दोनों चिह्न दोनों स्थितियों का अभिलक्षण कर देते हैं।

\textbf{19.} दुहरे एकांतरक में $b_i = a_i$ लीजिए: अंश $\prod_{i<j}(a_j -
a_i)^2 > 0$ है और हर $\prod_{i,j}(a_i + a_j) > 0$ (सभी प्रविष्टियाँ
धनात्मक): अतः सारणिक धनात्मक है। (आगे की भाषा में: नाभिक $\frac1{x+y}$
धन-निश्चित है।)

\textbf{20.} प्रश्न 2--3 से सारणिक $V(2,4,7)/(0!\,1!\,2!) =
\frac{(4-2)(7-2)(7-4)}{2} = \frac{30}{2} = 15$ के बराबर है। सीधे देखें तो
आव्यूह है
\[
\begin{pmatrix}
1 & 2 & 1\\
1 & 4 & 6\\
1 & 7 & 21
\end{pmatrix},
\qquad
\det = (84 - 42) - 2(21 - 6) + (7 - 4) = 42 - 30 + 3 = 15 .
\]

\textbf{21.} मान लीजिए अनुक्रम के रूप में $\sum_i c_i\,(\lambda_i^{\,k})_{k}
= 0$। $k = 0, 1, \dots, n-1$ पढ़ लेने पर $W^{\mathsf T}c = 0$ मिलता है, जहाँ
$W = (\lambda_i^{\,j-1})$ व्युत्क्रमणीय है ($\det = V \neq 0$, भिन्न
$\lambda_i$): अतः $c = 0$। इस प्रकार गुणोत्तर अनुक्रम \omterm{def:b1:vspaces:free}{स्वतंत्र} हैं।

\textbf{22.} $a = b = (1, 2, 3)$: अंश $\bigl[(2-1)(3-1) (3-2)\bigr]^2 = 4$;
हर $\prod_{i,j}(i + j) = (2\cdot3\cdot4)(3\cdot4\cdot5)(4\cdot5\cdot6) =
24\cdot60\cdot120 = 172800$। अतः $\det\bigl(\frac1{i+j}\bigr) =
\frac{4}{172800} = \frac1{43200}$।

\textbf{23.} यदि $x_i = x_j$, तो दोनों चरों की अदला-बदली उस बिंदु को अचल
रखती है पर $F$ का चिह्न बदल देनी ही चाहिए: $F = -F$, अतः वहाँ $F = 0$।
\omterm{def:b1:arith:divides}{विभाज्यता}: $F$ को अकेले चर $x_n$ में एक \omterm{def:b1:poly:def}{बहुपद} मानिए जिसके गुणांक शेष चरों
में हैं; वह $n - 1$ ``मानों'' $x_1, \dots, x_{n-1}$ पर शून्य होता है, अतः
बार-बार गुणनखंडन (\cref{thm:b1:poly:factor}) से $F = \prod_{i<n}(x_n -
x_i)\cdot G$ मिलता है, जहाँ $G$ \omterm{def:b1:poly:def}{बहुपद} है। पूर्व-गुणक दो सूचकांकों $i, j < n$
की अदला-बदली के अंतर्गत अपरिवर्तित है, अतः $G$ $x_1, \dots, x_{n-1}$ में
एकांतरक है, और आगमन काम पूरा कर देता है: $\prod_{i<j}(x_j - x_i)$ $F$ को
\omterm{def:b1:arith:divides}{विभाजित करता है}।

\textbf{24.} $D = \det(x_i^{\,j-1})$ $x_i$ में एक \omterm{def:b1:poly:def}{बहुपद} है; दो चरों की
अदला-बदली दो पंक्तियों की अदला-बदली है, अतः $D$ एकांतरक है, और प्रश्न 23 से
किसी \omterm{def:b1:poly:def}{बहुपद} $c$ के लिए $D = c\,\prod_{i<j}(x_j - x_i)$। कुल घातें: $D$ की घात
$\leq 0 + 1 + \dots + (n-1) = \binom n2$ है, और गुणनफल की घात ठीक $\binom
n2$: अतः $c$ एक अचर है। एकपदी $x_2\,x_3^2\cdots x_n^{\,n-1}$ का गुणांक $D$
में $1$ है (विकर्ण गुणनफल) और गुणनफल में $1$ (हर गुणनखंड में बड़े सूचकांक
वाला चर चुनिए): अतः $c = 1$, और वांडरमोंड सूत्र बिना किसी आगमन के निकल आता
है।

\textbf{25.} (क) एकांतरकता --- अर्थात् वह अभिगृहीत कि ``दो बराबर स्तंभ
सारणिक को मार देते हैं'' --- ही यंत्र है: उसी ने इस समस्या का हर गुणनखंड
$(x_j - x_i)$, $(a_j - a_i)$, $(b_j - b_i)$ पैदा किया। (ख) सर्वसमिका $\det S
= V^2$ व्यक्तिगत रूप से सम्मिश्र और अगम्य मूलों के स्थान पर उनके घात-योग रख
देती है, जो गुणांकों में वास्तविक \omterm{def:b1:poly:def}{बहुपद} हैं, अतः भिन्नता किसी संगणनीय
वास्तविक संख्या का चिह्न बन जाती है। (ग) \omterm{ex:b1:det:vandermonde}{वांडरमोंड सारणिक} बहुपदीय अंतर्वेशन
पर शासन करता है, और \omterm{pb:b1:det:1}{कोशी सारणिक} आंशिक भिन्नों तथा प्रतिचयित परिमेय फलनों पर
(जिसकी सबसे प्रसिद्ध विशेष स्थिति \omterm{pb:b1:det:1}{हिल्बर्ट आव्यूह} है)। (घ) कोई भी एकांतरक
\omterm{def:b1:poly:def}{बहुपद} $\prod_{i<j}(x_j - x_i)$ से विभाज्य होता है, और फिर घात की गिनती ऐसे
\omterm{def:b1:poly:def}{बहुपद} को किसी अचर तक कील देती है --- और इसीलिए यह गुणनफल वहीं-वहीं फिर से
आता रहता है जहाँ कोई सारणिक संपातों पर शून्य होता है। भाग III की प्रमेय
\emph{कोशी का दुहरा एकांतरक} है।
\end{solution}
